\cleardoublepage
\chapter{PENUTUP}
\label{chap:penutup}

\section{Kesimpulan}
Berdasarkan rangkaian tahapan implementasi kerangka kerja, pengujian skenario studi ablasi, serta analisis pembahasan terhadap arsitektur pemodelan sekuensial untuk estimasi konsentrasi PM\textsubscript{10} di DKI Jakarta, diperoleh kesimpulan yang menjawab seluruh pertanyaan penelitian secara terstruktur sebagai berikut:

\begin{enumerate}
    \item {Rancangan Model Estimasi Harian $PM_{10}$ ($t+1$):} Model prediksi harian konsentrasi PM\textsubscript{10} untuk proyeksi satu hari ke depan ($t+1$) berhasil dirancang menggunakan basis data deret waktu historis Indeks Standar Pencemar Udara (ISPU) DKI Jakarta periode tahun 2010--2025 dengan resolusi harian. Model ini dikonstruksi secara mandiri untuk masing-masing dari lima stasiun pemantau guna menangkap karakteristik lokal, dengan mengandalkan horizon waktu dinamis di mana jendela data riwayat bergerak maju secara kronologis untuk memproyeksikan nilai target berikutnya.
    \item {Mekanisme Integrasi Sekuensial Model Hibrida:} Integrasi antara algoritma berbasis pohon keputusan dan model statistik dilakukan melalui arsitektur hibrida sekuensial dua fase. Pada fase primer, algoritma \textit{Random Forest Regressor} digunakan untuk memproses matriks fitur komposit guna memetakan tren fluktuasi non-linier data multivariat polutan. Selanjutnya, fase sekunder memanfaatkan model statistik ARIMA untuk mengolah deret sisaan galat (\textit{residual}) linier dari fase pertama yang terbukti masih menyimpan informasi autokorelasi dan dependensi temporal runtun waktu. Penjumlahan hasil prediksi dari kedua fase ini menghasilkan nilai estimasi akhir dengan autokorelasi galat yang terkendali.
    \item {Pengaruh Penerapan Rekayasa Fitur terhadap Galat Estimasi:} Hasil eksperimen melalui skema studi ablasi secara empiris menunjukkan bahwa pengayaan dimensi masukan lewat \textit{feature engineering} menghasilkan nilai kesalahan prediksi yang lebih rendah. Penambahan informasi dependensi temporal berupa fitur jeda waktu (\texttt{pm10\_lag1}), statistik berjalan (\texttt{pm10\_ma3}, \texttt{pm10\_ma7}, \texttt{pm10\_std7}), dan fitur interaksi kimiawi ($CO \times O_3$) berhasil menurunkan metrik galat \textit{Root Mean Square Error} (RMSE) set pengujian secara bertahap dari 11,45 pada Model A (fitur polutan mentah murni) menjadi 7,12 pada Model C (model penuh usulan). Penurunan nilai kesalahan ini juga disertai dengan capaian nilai Koefisien Determinasi ($R^2$) yang menyentuh angka 0,89 pada model penuh.
    \item {Analisis Transparansi Perilaku Model Menggunakan SHAP:} Penerapan pendekatan \textit{Explainable AI} (XAI) melalui metode \textit{TreeSHAP} pada hasil prediksi fase primer terbukti valid untuk menggambarkan kontribusi marginal aditif dari setiap prediktor secara terukur. Visualisasi melalui grafik sebaran global (\textit{beeswarm plot}) menunjukkan bahwa riwayat jangka pendek hari sebelumnya (\texttt{pm10\_lag1}) memiliki bobot atribusi terbesar dalam keputusan model, sedangkan grafik ketergantungan (\textit{dependence plot}) memetakan arah hubungan nonlinear antar-polutan prekursor sekunder yang secara ilmiah selaras dengan karakteristik domain sains atmosfer wilayah DKI Jakarta.
\end{enumerate}

\section{Saran}
Berdasarkan batasan operasional dan kendala teknis yang diidentifikasi selama pelaksanaan eksperimen Tugas Akhir ini, beberapa saran yang dapat diajukan untuk riset atau pengembangan selanjutnya adalah sebagai berikut:

\begin{enumerate}
    \item {Integrasi Variabel Meteorologi Makro:} Penelitian selanjutnya dapat mempertimbangkan untuk memasukkan variabel eksogen berupa parameter meteorologi luar, seperti intensitas curah hujan, kecepatan angin, arah sirkulasi angin, serta tingkat kelembapan udara. Penambahan parameter fisik ini ditujukan untuk mengurangi sisa nilai galat estimasi model saat memproses fenomena peluruhan polutan oleh air hujan (\textit{wet deposition}) maupun perpindahan massa partikulat antarwilayah secara spasial.
    \item {Pengembangan Horizon Peramalan \textit{Multi-Step Forecasting}:} Mengingat cakupan proyeksi dalam penelitian ini masih terbatas pada horizon satu hari ke depan (\textit{single-step forecasting}), pengembangan berikutnya dapat diarahkan untuk memproyeksikan rentang waktu yang lebih panjang (misalnya 3 hingga 7 hari ke depan). Skenario jangka menengah tersebut dapat dikombinasikan atau dibandingkan dengan arsitektur pemrosesan sekuensial lainnya seperti \textit{Long Short-Term Memory} (LSTM) untuk mengamati perbandingan stabilitas nilai galatnya.
    \item {Automasi Sistem Alur Pemrosesan Data:} Model estimasi harian yang telah diuji dapat dikembangkan lebih lanjut menjadi komponen fungsional pendukung dasbor pemantauan milik instansi terkait. Proses pengembangan tersebut memerlukan penerapan mekanisme otomasi penarikan data dari sensor secara berkala setiap hari, serta perancangan sistem penanganan gangguan transmisi data (\textit{data drop assurance}) menggunakan arsitektur komputasi awan yang sinkron.
\end{enumerate}